% !TEX root =  ../ProjectSpecification.tex

\section{Resources \& Risks}
\label{sec:ResourcesAndRisks}

\subsection{Resources}
\large{\raggedright
Preliminary research has included studying the fundamentals of QC using sections 1 and 2 of the Quantum Information and Computation lecture notes by Richard Jozsa \cite{Jozsa19}, alongside completing the first three modules of the \href{https://pennylane.ai/codebook#modules}{\color{blue} PennyLane Codebook}. Further materials consulted throughout the project will be referenced in subsequent reports.\\
~\\
\href{https://github.com/Freddie-Audley/3rd-Year-Dissertation}{\color{blue} GitHub} is being used for version control, providing backups in the event of data loss and a record of project development through its commit history.\\
~\\
The available hardware includes a PC and a laptop, although neither is suited to QML training beyond the implementation of the pipeline. The DCS Batch Compute system\footnotemark[1] will be used to train the larger scale models from which data will be collected.
\par
}

\subsection{Risks}
\large{\raggedright
Small scale experiments may not be sufficiently expressive to support interpretation. Access to the DCS Batch Compute\footnotemark[1] system can help mitigate this limitation by providing additional computational resources for larger experiments, although the feasible scale will depend on the complexity of the circuits and the simulation method used.\\
~\\
Given the specialised nature of the project, additional time has been allocated for preliminary research. The timetable also incorporates buffers to accommodate external academic commitments and other unforeseen circumstances.
\par
}

\footnotetext[1]{\hangindent=1.8em \url{https://warwick.ac.uk/fac/sci/dcs/admissions/postgraduatetaughtdegrees/facilities_research/} "High-performance computing"}